\documentclass[10pt,a4paper]{amsart}
\usepackage[margin=19mm]{geometry}
\usepackage{amsmath,amssymb,mathtools}
\usepackage[scr=rsfs,cal=euler]{mathalfa}
\usepackage{xfrac}
\usepackage[unicode,hidelinks,bookmarks=false]{hyperref}
\newcommand{\ie}{i.e.\ }
\newcommand{\cf}{cf.\ }
\newcommand{\resp}{respectively\ }
\newcommand{\Subsection}{\subsection}
\providecommand{\nobreakdash}{\nobreak\hbox{-}}
\setlength{\emergencystretch}{2em}
\multlinegap0pt
\usepackage[shortlabels]{enumitem}
\usepackage{amssymb}
\usepackage{tikz}
\newtheorem{lemma}{Lemma}
\newtheorem{theorem}{Theorem}
\newtheorem{corollary}{Corollary}
\newcommand{\M}{\mathrm M_\infty(K)}
\newcommand{\Span}{\operatorname{span}_K}
\newcommand{\T}{\mathscr T_K}
\title{Lie Ideals and Lie Derivations of the Algebraic Toeplitz Algebra}
\author{Nguyen Huynh Thao Nhi, Huynh Viet Khanh}
\date{Source: arXiv:2609.04745v1 (CC BY 4.0)}
\begin{document}
\maketitle

\noindent\emph{Source and licence.} Lie Ideals and Lie Derivations of the Algebraic Toeplitz Algebra. Version arXiv:2609.04745v1 (CC BY 4.0), distributed by arXiv under the Creative Commons Attribution 4.0 International licence, http://creativecommons.org/licenses/by/4.0/, as recorded in the arXiv licence metadata for this exact version and shown on the abstract page https://arxiv.org/abs/2609.04745v1. Full licence notice: \texttt{licenses/arxiv-2609.04745-CC-BY-4.0.txt}.\par\smallskip

\noindent\textit{Editorial scope.} Counted unit: the article's corollary cor:graded-lie-ideals (source0.tex lines 497--505). The article's complete original proof of it is delivered with it (source0.tex lines 507--605, one \texttt{proof} environment of 99 lines). No mathematics is written, repaired or re-derived: the statement and the proof are the article's own lines, in the article's own notation, and no retained line is substituted or re-flowed.  Retained uncounted as the source-local context the proof needs: lines 137--139; lines 147--156; lines 164--166; lines 171--173; lines 410--417.  Nothing was omitted from any retained span, so the package carries no omission record.  No retained line contains a citation, so the wrapper carries no bibliography and no bibliography entry.  No retained line contains a figure environment, an image-inclusion command or drawing code, so no image is delivered and the package declares no asset dependency.  Editorial formatting only, in addition: an AMS \texttt{amsart} preamble that restates the article's own declarations and macros used by the retained lines, namely the environments \texttt{lemma}, \texttt{theorem}, \texttt{corollary} on separate counters, together with the standard packages those lines need.  The retained environments print in this document as Lemma 1, Theorem 1, Corollary 1 in order of appearance, whereas the article prints the same items with the numbering of its own full document.  The complete native source of the article as distributed in the arXiv e-print for this exact version is bundled unchanged as \texttt{sources/209385/source0.tex}.
\begin{lemma}\label{lem:matrix-units}
Let $I(w)$ be the ideal of $\T$ genenerated by $w$. For all $i,j,k,l\geq1$, $F_{ij}F_{kl}=\delta_{jk}F_{il}$. Consequently $F:=\bigoplus_{i,j\geq1}KF_{ij}$ is a two-sided ideal of $\T$, $F\cong\M$, and $F=I(w)$.
\end{lemma}

\begin{lemma}\label{lem:path-basis}
The set
$$
        \mathcal B=\{v\}\cup\{e^n:n\geq1\}\cup\{(e^*)^n:n\geq1\}\cup\{F_{ij}:i,j\geq1\}
$$
is a $K$-basis of $\T$. Consequently, as a vector space, we have
$$
        \T=K1\oplus\Bigl(\bigoplus_{n\geq1}Ke^n\Bigr)\oplus\Bigl(\bigoplus_{n\geq1}K(e^*)^n\Bigr)\oplus F.
$$
\end{lemma}

\begin{lemma}\label{lem:quotient}
Let $\rho:\T\to K[t,t^{-1}]$ be the homomorphism determined by $\rho(v)=1$, $\rho(e)=t$, $\rho(e^*)=t^{-1}$, and $\rho(w)=\rho(f)=\rho(f^*)=0$. Then $\ker\rho=F$, and hence $\T/F\cong K[t,t^{-1}]$.
\end{lemma}

\begin{lemma}\label{lem:s-trace-zero}
One has $F=KF_{11}\oplus\mathfrak s$ as vector spaces.
\end{lemma}

\begin{theorem}\label{thm:classification}
The Lie ideals of $\T^-$ are precisely the following:
\begin{enumerate}[label=\rm(\roman*)]
\item $0$ and $K1$;
\item $\mathfrak s+U$, where $U\leq P_K\oplus Kw$ and $w\notin U$;
\item $\rho^{-1}(W)$, where $W\leq K[t,t^{-1}]$ is a $K$-subspace.
\end{enumerate}
\end{theorem}

\begin{corollary}[Graded Lie ideals]\label{cor:graded-lie-ideals}
Give $\T$ its natural $\mathbb Z$-grading. For $n\in\mathbb Z$, put $F_n=\Span\{F_{ij}:i-j=n\}$. Then
$$
{\T}_n=K e^n\oplus F_n\quad\text{if }n>0,\qquad
{\T}_0=K1\oplus F_0,\qquad
{\T}_n=K(e^*)^{-n}\oplus F_n\quad\text{if }n<0.
$$
Moreover, $F=\bigoplus_{n\in\mathbb Z}F_n$, and $\mathfrak s$ is graded, with $\mathfrak s_n=F_n$ for $n\neq0$ and $\mathfrak s_0=\left\{\sum_i a_iF_{ii}:\sum_i a_i=0\right\}$. The graded Lie ideals of $\T^-$ are precisely $0$, $K1$, the spaces $\mathfrak s+U$, where $U$ is a graded subspace of $P_K\oplus Kw$ and $w\notin U$, and the spaces $\rho^{-1}(W)$, where $W$ is a graded subspace of $K[t,t^{-1}]$.
\end{corollary}

\begin{proof}
By Lemma~\ref{lem:path-basis},
$$
\mathcal B=\{v\}\cup\{e^r:r\geq1\}\cup\{(e^*)^r:r\geq1\}\cup\{F_{ij}:i,j\geq1\}
$$
is a $K$-basis of $\T$. Each element of $\mathcal B$ is homogeneous. More precisely,
$$
\deg v=0,\qquad \deg e^r=r,\qquad\text{and}\qquad \deg (e^*)^r=-r.
$$
Recall that $\alpha_1=w$, $\alpha_i=e^{i-2}f$ for $i\geq2$, and $F_{ij}=\alpha_i\alpha_j^*$. Since $\deg w=0$ and
$$
\deg(e^{i-2}f)=i-1\qquad\text{for }i\geq2,
$$
one has $\deg\alpha_i=i-1$ for all $i\geq1$. It follows that
$$
\deg F_{ij}=\deg\alpha_i-\deg\alpha_j=(i-1)-(j-1)=i-j.
$$

Fix $n>0$. The elements of $\mathcal B$ of degree $n$ are $e^n$ and the matrix units $F_{ij}$ satisfying $i-j=n$. Hence ${\T}_n=Ke^n\oplus F_n$. The basis elements of degree $0$ are $v$ and the diagonal matrix units $F_{ii}$. Since $w=F_{11}\in F_0$ and $v=1-w$, one has $Kv+F_0=K1+F_0$. This sum is direct because $\rho(1)=1$ and $\rho(F_0)=0$ by Lemma~\ref{lem:quotient}. Thus ${\T}_0=K1\oplus F_0$. Finally, if $n<0$, the elements of $\mathcal B$ of degree $n$ are $(e^*)^{-n}$ and the matrix units $F_{ij}$ satisfying $i-j=n$. Therefore ${\T}_n=K(e^*)^{-n}\oplus F_n$.

Since the $F_{ij}$ form a $K$-basis of $F$ by Lemma~\ref{lem:matrix-units}, and each $F_{ij}$ belongs to exactly one $F_n$, we obtain $F=\bigoplus_{n\in\mathbb Z}F_n$. Put $\mathfrak s_n=\mathfrak s\cap{\T}_n$. By Lemma~\ref{lem:s-trace-zero},
$$
\mathfrak s=\operatorname{span}_K\{F_{ij}:i\neq j\}
+\operatorname{span}_K\{F_{ii}-F_{jj}:i,j\geq1\}.
$$
If $n\neq0$ and $i-j=n$, then $i\neq j$, so $F_{ij}\in\mathfrak s$. Hence $F_n\subseteq\mathfrak s_n$. Conversely, $\mathfrak s\subseteq F$ gives $\mathfrak s_n\subseteq F_n$, and therefore $\mathfrak s_n=F_n$ for $n\neq0$.
Since $F_0=\bigoplus_{i\geq1}KF_{ii}$, it follows from  Lemma~\ref{lem:s-trace-zero} that
$$
\mathfrak s_0
=\ker\bigl(\operatorname{Tr}|_{F_0}\bigr)
=\Big\{\sum_i a_iF_{ii}:\text{only finitely many }a_i\neq0
\text{ and }\sum_i a_i=0\Big\}.
$$
Thus $\mathfrak s=\mathfrak s_0\oplus\bigoplus_{n\neq0}F_n$, so $\mathfrak s$ is graded.

We now determine which Lie ideals in Theorem~\ref{thm:classification} are graded. The spaces $0$ and $K1$ are graded. The space $P_K\oplus Kw$ is also graded. If $\operatorname{char}K=0$, then
$$
P_K\oplus Kw=K1\oplus Kw\subseteq{\T}_0.
$$
If $\operatorname{char}K=\ell>0$, then
$$
P_K\oplus Kw
=(K1\oplus Kw)
\oplus\bigoplus_{m\geq1}Ke^{m\ell}
\oplus\bigoplus_{m\geq1}K(e^*)^{m\ell},
$$
where the displayed summands have degrees $0$, $m\ell$, and $-m\ell$, respectively.

Let $U\leq P_K\oplus Kw$ with $w\notin U$. If $U$ is graded, then $\mathfrak s+U$ is graded because both $\mathfrak s$ and $U$ are graded. Conversely, suppose that $\mathfrak s+U$ is graded. The quotient map $\pi:\T\to\T/\mathfrak s$ is graded because $\mathfrak s$ is graded. Moreover, as established immediately before Theorem~\ref{thm:classification}, the restriction
$$
\pi|_{P_K\oplus Kw}:P_K\oplus Kw\longrightarrow Z(\T/\mathfrak s)
$$
is injective. Hence
$$
\mathfrak s\cap(P_K\oplus Kw)=0
\qquad\text{and}\qquad
(\mathfrak s+U)\cap(P_K\oplus Kw)=U.
$$
Let $u\in U$, and write its homogeneous decomposition in the graded space $P_K\oplus Kw$ as
$$
u=\sum_{n\in\mathbb Z}u_n,\qquad\text{where } u_n\in(P_K\oplus Kw)\cap{\T}_n.
$$
Since $u\in\mathfrak s+U$ and $\mathfrak s+U$ is graded, every $u_n$ belongs to $\mathfrak s+U$. The preceding intersection identity then gives $u_n\in U$ for every $n$. Thus $U$ is graded. We have proved that
$$
\mathfrak s+U\text{ is graded}
\quad\text{if and only if}\quad
U\text{ is graded}.
$$

Next let $W\leq K[t,t^{-1}]$. By Lemma~\ref{lem:quotient} and the formulas for the homogeneous components of $\T$ proved above,
$$
\rho({\T}_n)\subseteq Kt^n\qquad\text{for all }n\in\mathbb Z.
$$
Thus $\rho$ is a graded homomorphism, where $K[t,t^{-1}]$ has its standard grading $K[t,t^{-1}]=\bigoplus_{n\in\mathbb Z}Kt^n$. If $W$ is graded, then $\rho^{-1}(W)$ is graded. Indeed, let $x\in\rho^{-1}(W)$ and write
$$
x=\sum_n x_n,\qquad\text{where } x_n\in{\T}_n.
$$
The elements $\rho(x_n)\in Kt^n$ are the homogeneous components of $\rho(x)\in W$. Since $W$ is graded, $\rho(x_n)\in W$ for every $n$, and hence $x_n\in\rho^{-1}(W)$.

Conversely, suppose that $\rho^{-1}(W)$ is graded. Let $f\in W$. Since $\rho$ is surjective, choose $x\in\rho^{-1}(W)$ with $\rho(x)=f$, and write $x=\sum_nx_n$ with $x_n\in{\T}_n$. The gradedness of $\rho^{-1}(W)$ gives $x_n\in\rho^{-1}(W)$ for every $n$. Hence $\rho(x_n)\in W$, and these are precisely the homogeneous components of $f$. Therefore $W$ is graded. We have proved that
$$
\rho^{-1}(W)\text{ is graded}
\quad\text{if and only if}\quad
W\text{ is graded}.
$$

Theorem~\ref{thm:classification} now gives the asserted list of graded Lie ideals. Finally, since each homogeneous component $Kt^n$ of $K[t,t^{-1}]$ is one-dimensional, a subspace $W\leq K[t,t^{-1}]$ is graded if and only if
$$
W=\bigoplus_{n\in S}Kt^n
=\operatorname{span}_K\{t^n:n\in S\}
$$
for some subset $S\subseteq\mathbb Z$. Similarly, if $\operatorname{char}K=\ell>0$, a graded subspace of $P_K\oplus Kw$ has the form
$$
U=U_0
\oplus\bigoplus_{m\in S_+}Ke^{m\ell}
\oplus\bigoplus_{m\in S_-}K(e^*)^{m\ell},
$$
where $U_0\leq K1\oplus Kw$, $S_+,S_-\subseteq\mathbb N$, and the condition $w\notin U$ is equivalent to $w\notin U_0$.
\end{proof}

\end{document}
